\documentclass[../main.tex]{subfiles}
\begin{document}

% --------------------------------------------------
\section{note.sty}
% --------------------------------------------------

\texttt{note.sty}は, このテンプレート全体を読み込むための入口となるスタイルファイルである. 
このテンプレートをそのまま使う場合は, プリアンブルに
\begin{center}
	\verb|\usepackage{sty/note}|
\end{center}
と書けばよい. 各スタイルファイルを個別に読み込む必要はない. 

\texttt{note.sty}では, はじめに\texttt{core.sty}を読み込む. 
\texttt{core.sty}で色や基本パッケージを定義してから, 他のスタイルファイルがそれらを使うためである. 
その後, 文献, コード用の予約ファイル, レイアウト, 数式, リンク, 参照, 定理環境の順に読み込む. 
\begin{center}
	\texttt{core, bib, code, layout, math, hyperref, bookmark, ref, theorem}
\end{center}

リンクまわりについては, \texttt{hyperref.sty}と\texttt{bookmark.sty}を読み込んでいる. 
PDFのしおりは番号付きにし, 本文中のリンク枠は表示しない設定にしている. 

色, 見出し, 定理環境, 数式マクロ, 参照, 文献表示などの詳細は, それぞれ対応する節で説明する. 


\end{document}
